\newcommand{\sujet}[1]{\renewcommand{\sujet}{#1}}
\newcommand{\auteur}[1]{\renewcommand{\auteur}{#1}}
\newcommand{\encadrant}[1]{\renewcommand{\encadrant}{#1}}
\makeatletter
\renewcommand{\tableofcontents}{%
  \@starttoc{toc}%
}
\makeatother
\newcommand{\chapminitoc}{%
  \begingroup
  \etocsettocstyle{\small}{}
  \localtableofcontents
  \endgroup
}
\newcommand{\objectif}[1]{%
\vspace{.1cm}
\begin{singlespace}
\tikzstyle{titlebox}=[rectangle,inner sep=10pt,inner ysep=10pt,draw=curcolor,draw]%
\tikzstyle{title}=[fill=white]%
\bigskip\noindent\begin{tikzpicture}
\node[titlebox] (box){%
    \begin{minipage}{0.88\textwidth}
#1
    \end{minipage}
};
\node[title] at (box.north west) {\color{curcolor}  Abstract};
\end{tikzpicture}\bigskip%
\vspace{.1cm}
\chapminitoc
\end{singlespace}
\newpage
}

\newcommand{\myparagraph}[1]{\paragraph{#1}\mbox{} \vskip .5\baselineskip \par}

\newcommand{\T}[1]{T\textsubscript{#1}}


%%% Macro pour la mise en forme de la liste des publiations
\newcommand{\Myprod}[3]{ % sans DOI
\begin{singlespace}
  \noindent {#1}. "\textit{{#2}}". {#3}. \par\vspace{0.6em}
\end{singlespace}}

\newcommand{\MyprodwithDOI}[4]{ % avec DOI
\begin{singlespace}
  \noindent {#1}. "\textit{{#2}}". {#3}.\par
 \noindent {DOI: \href{https://doi.org/#4}{#4}} \par\vspace{0.6em}
\end{singlespace}}

% Chapitre sans mise en forme particulière (remerciements, conclusion
\newcommand{\addchapnonumber}[1]{
\phantomsection
\addtocounter{chapter}{1}
\chapter*{#1}
\addcontentsline{toc}{chapter}{#1}
\markboth{#1}{#1}
\setcounter{section}{1}
}

%%% Macro pour mise en forme de l'abstract et résumé avec mots clés
\newcommand{\AddResumeAbstract}{
\chapter{Résumé}

\thefrabstract
\vskip 2em \noindent\makebox[\linewidth]{\rule{.5\linewidth}{0.4pt}}
\vfill
\noindent \textbf{Mots clés :}
\thefrkeywords

\chapter{Abstract}

\theenabstract
\vskip 2em \noindent\makebox[\linewidth]{\rule{.5\linewidth}{0.4pt}}
\vfill
\noindent\textbf{Keywords :}
\theenkeywords
}

\makeatletter
\newcommand{\mytag}[2]{%
  \text{#1}%
  \@bsphack
  \begingroup
    \@onelevel@sanitize\@currentlabelname
    \edef\@currentlabelname{%
      \expandafter\strip@period\@currentlabelname\relax.\relax\@@@%
    }%
    \protected@write\@auxout{}{%
      \string\newlabel{#2}{%
        {#1}%
        {\thepage}%
        {\@currentlabelname}%
        {\@currentHref}{}%
      }%
    }%
  \endgroup
  \@esphack
}
\makeatother
\numberwithin{equation}{section}
\renewcommand{\theequation}{\arabic{section}.\arabic{equation}}
\usepackage[shortlabels]{enumitem}
\newlist{mathlist}{enumerate}{5}
\setlist[mathlist]{label={\textup{(\roman*)}}}
\theoremstyle{plain}
\newtheorem{thm}{Theorem}[section]
\newtheorem{propdef}[thm]{Definition/Proposition}
\newtheorem{prop}[thm]{Proposition}
\newtheorem{lem}[thm]{Lemma}
\newtheorem{cor}[thm]{Corollary}
\theoremstyle{definition}
\newtheorem{assumption}{Assumption}
\renewcommand{\theassumption}{\Roman{assumption}}
\newtheorem{algo}[thm]{Algorithm}
\newtheorem{defn}[thm]{Definition}
\newtheorem{notation}[thm]{Notation}
\theoremstyle{remark}
\newtheorem{rem}[thm]{Remark}
\renewcommand{\qedsymbol}{\ensuremath{\blacksquare}}
\renewcommand{\thethm}{\arabic{section}.\arabic{thm}}
\allowdisplaybreaks[1]

\renewcommand{\arraystretch}{1.5}
\newcommand{\norm}[2][]{#1\lVert #2 #1\rVert}
\newcommand{\abs}[2][]{#1\lvert #2 #1\rvert}
\newcommand{\floor}[2][]{#1\lfloor #2 #1\rfloor}
\renewcommand{\rho}{\varrho}
\renewcommand{\epsilon}{\varepsilon}
\newcommand{\ppac}[1]{\llbracket#1\rrbracket}
\DeclareMathOperator{\R}{\mathbb{R}}


\usepackage{setspace}
\makeatletter
\newcommand{\oset}[3][0.1ex]{%
  \mathrel{\mathop{#3}\limits^{
    \vbox to#1{\kern-2\ex@
    \hbox{$\scriptstyle#2$}\vss}}}}
\makeatother
\makeatletter
\newcommand*{\transpose}{%
  {\mathpalette\@transpose{}}%
}
\newcommand*{\@transpose}[2]{%
  % #1: math style
  % #2: unused
  \raisebox{\depth}{$\m@th#1\intercal$}%
}
\makeatother


\renewcommand{\le}{\leqslant}
\renewcommand{\ge}{\geqslant}
\renewcommand{\leq}{\leqslant}
\renewcommand{\geq}{\geqslant}
\renewcommand{\baselinestretch}{1.05}
\renewcommand{\arraystretch}{1.2}
\renewcommand{\epsilon}{\varepsilon}
\newcommand{\eps}{\varepsilon}
\renewcommand{\rho}{\varrho}
\DeclareMathOperator{\Var}{Var}
\newcommand{\1}{\mathds{1}}
\newcommand{\X}{\mathcal{X}}
\newcommand{\ie}{\textit{i.e.}\ }
\newcommand{\eg}{\textit{e.g.}\ }
\newcommand{\viz}{\textit{viz.}\ }
\newcommand{\cf}{\textit{cf.}\ }
\newcommand{\iid}{i.i.d.\ }
\DeclareMathOperator{\Poi}{\mathrm{Poi}}
\newcommand{\indep}{\raisebox{0.05em}{\rotatebox[origin=c]{90}{$\models$}}}
\newcommand{\anc}[1]{\overset\leftarrow{#1}}
\DeclareMathOperator{\E}{\mathbb{E}}
\DeclareMathOperator{\N}{\mathbb{N}}
\DeclareMathOperator{\Z}{\mathbb{Z}}
\DeclareMathOperator{\Q}{\mathbb{Q}}
\DeclareMathOperator{\Prob}{\mathbb{P}}
\DeclareMathOperator{\Cte}{C^{\textit{st}}}
\renewcommand{\figurename}{Fig.}
\newcommand{\ma}[1]{\textcolor{violet}{\ \textbf{MA:} #1}}
\newcommand{\jj}[1]{\textcolor{brown}{\textbf{JJD:} #1}}
\newcommand{\jjn}[1]{\textcolor{brown}{#1}}

\newcommand{\ffr}[1]{\textcolor{orange}{\textbf{F}: #1}}

\newcommand{\es}[1]{\textcolor{red}{ES: #1}}

\newcommand{\esn}[1]{\textcolor{red}{#1}}